\chapter*{Introduzione al corso}
\addcontentsline{toc}{chapter}{Introduzione al corso}
Il corso di Organizzazione e gestione per lo start-up aziendale si propone di esplorare il mondo delle imprese e il loro avvio. Il corso si focalizza sulle fasi cruciali della creazione di un'azienda, le tecniche e gli strumenti necessari per trasformare un'idea di business in un'impresa funzionante.

Il corso è strutturato in tre macro-aree:
\begin{enumerate}
    \item \textit{Impresa e organizzazione:} una panoramica sui concetti di base dell’impresa, i suoi attori e le fasi del ciclo di vita.
    \item \textit{Start-up:} si parte dall’idea imprenditoriale per giungere alla formalizzazione del modello di business, passando per analisi di fattibilità di mercato ed economico-finanziaria.
    \item \textit{Business plan:} copre aspetti come la value proposition, le ipotesi di base e le proiezioni economiche future.
\end{enumerate}
Il corso ha tre obiettivi principali:
\begin{enumerate}
    \item Fornire una visione "idonea" e moderna dell’impresa nel contesto competitivo attuale, analizzando le relazioni interne ed esterne all’impresa e adottando un approccio sistemico.
    \item Fornire degli strumenti necessari per sviluppare un’idea di business, includendo l’utilizzo del Business Model Canvas e la progettazione della value proposition.
    \item Insegnare come redigere un business plan, attraverso l’analisi di mercato, l’individuazione dei ricavi, la stima dei costi operativi e la realizzazione di proiezioni economico-finanziarie.
\end{enumerate}
Alcuni concetti chiave del corso includono: organizzazione dell’impresa, modello di business, conti economici, SWOT analysis, e flussi di cassa. Questi termini rappresentano gli strumenti fondamentali per comprendere e gestire l’avvio di un’impresa in un contesto di mercato globale e competitivo.

Si utilizzeranno gli strumenti per esplicitare un'idea di business, sviluppandola fino alla sua formalizzazione in un modello di business operativo. Inoltre, si acquisirà la capacità di comprendere e redigere un business plan, un documento strategico fondamentale per la gestione e la crescita dell'impresa.